\section{Related work and discussion}\label{sec:related}

\paragraph{Viability and invariance.}
The kernel of Definition~\ref{def:game} is the discrete-time viability kernel of Aubin \cite{Aubin1991} and the maximal controlled-invariant set of Bertsekas and Rhodes \cite{BertsekasRhodes1971} and Blanchini \cite{Blanchini1999}, computed on the finite lattice of beliefs instead of a continuous state space; the greatest-fixed-point iteration is the finite counterpart of the viability algorithm of Saint-Pierre \cite{SaintPierre1994}, and the exactness it affords is what the approximation methods of \cite{Deffuant2007} trade for scale. The audit's contribution is not the algorithm but the systematic variation of the information and protocol axes with exact certificates for each inclusion and each strictness.

\paragraph{Games and processes with imperfect information.}
Beliefs as sets of states and fibrewise control are the knowledge-based subset construction of Reif \cite{Reif1984} and the antichain algorithms for safety games with imperfect information of Chatterjee, Doyen, Henzinger and Raskin \cite{CDHR2007}; the down-closure of Lemma~\ref{lem:basic}(b) is the reason antichains suffice there, and the maximal beliefs of Table~\ref{tab:kernels} are these antichains. The qualitative (almost-sure or sure) analysis of partially observable Markov decision processes \cite{CDH2010} and the belief-state formulation of POMDPs \cite{SmallwoodSondik1973,Kaelbling1998} ask the same question with probabilities; the sure-safety question studied here is the one for which exact finite answers exist, and its probabilistic refinement on related systems is the subject of \cite{Abaee2026prob}. Observability and co-observability of partitions in supervisory control \cite{LinWonham1988,CassandrasLafortune2008} are the discrete-event precedent for Section~\ref{sec:monitor}: there a partition is adequate when a supervisor can enforce a specification through it, here when the belief game is won through it, and Theorem~\ref{thm:dynamic} shows the dynamic question collapsing to a one-step graph condition on the viable set.

\paragraph{Information structures and teams.}
That the register convention changes the kernel is a small instance of Witsenhausen's observation that the information structure, not only the information, determines a control problem \cite{Witsenhausen1971,Ho1980}; the private register is a classical (nested) information pattern, the shared ledger a non-classical one in which a public record is written by several threads. The decentralised variant is a team problem in the sense of Radner \cite{Radner1962}, with the hardness results of \cite{TsitsiklisAthans1985,Bernstein2002} explaining why the factoring of Theorem~\ref{thm:factor}---which makes the problem tractable---depends on unanimity being statewise.

\paragraph{Max-plus, duality, Helly, partitions.}
The timing recursion of Proposition~\ref{prop:maxplus} is linear in the max-plus semiring \cite{Baccelli1992,Heidergott2006}; the continuous certificates are instances of Farkas' lemma and von Neumann's minimax theorem \cite{Farkas1902,vonNeumann1928}, and the sparsity bound is Helly's theorem \cite{Helly1923,DanzerGrunbaumKlee1963}; the count \eqref{eq:count} is a restricted Bell number in the sense of set-partition enumeration \cite{Stanley2012}. The exact-reproducibility standard follows the verified-artifact practice of the ARCH competitions \cite{Althoff2020,Geretti2020}, and the mechanised layer uses Lean~4 \cite{MouraUllrich2021}.

\paragraph{Limits.}
The worked system has two instruments, and Lemma~\ref{lem:alt} uses this: with three or more instruments the no-repeat protocol is not alternation, the closed forms of Theorem~\ref{thm:closed} do not apply, and the Helly number of the instrument family rises above $2$, so the graph criterion of Theorem~\ref{thm:adequacy} becomes necessary only. Theorem~\ref{thm:dynamic} is general but concerns the viable set as a whole; the kernel below it is a dynamic object with no one-step description, as the $74$ distinct kernels of the census show. The ledger convention is an auditing convention---an adversarial recording of one of several issued instruments---and the private register is the physically natural one; the article keeps both because the earlier audits are stated in the former and because the comparison is itself informative. Finally, the counts are those of one small family; the scaling to caps $4$ and $5$ shows the qualitative picture stable, but no asymptotic claim is made.
